% ============================================================
\section{Search and Discovery in Open-Ended Design Spaces}
\label{sec:ldm}
% ============================================================

In this work, we aim to discover a design or hypothesis $x\in\X$ that achieves the highest
reward $R(x)$, i.e.\ to find or closely approximate
$x^\star\in\argmax_{x\in\X}R(x)$. This setting presents two intrinsic
difficulties. First, $R$ is a black-box function that is expensive to evaluate
and potentially noisy. Second, the hypothesis space $\X$ is typically open-ended: vast,
combinatorial, structured, or not even fully specified in advance, as in the
design of molecules, biological sequences, computer programs, and experimental
workflows.

This differs from classical numerical optimisation or combinatorial search,
which generally assumes a fixed and explicitly specified domain together with
a gradient, numerical oracle, or known set of valid search operations. In
scientific discovery, the full hypothesis space may not be enumerable, and the
search may need to formulate hypotheses that lie beyond its current
representation or reach.

To reason about what is currently known, we introduce
\emph{a probabilistic surrogate}, i.e.\ a posterior
$p(R\mid\mathcal{D}_t)$ over the reward surface $R$ conditioned on the dataset
$\mathcal{D}_t=\{(x_i,r_i)\}_{i=1}^{t-1}$ of previously evaluated designs and
their noisy feedback $r_i\approx R(x_i)$. Its predictive belief is usually
summarised by posterior moments $(\mu_t(x),\sigma_t(x))$, representing the
expected reward and predictive uncertainty on the surrogate's modelling
support. The surrogate therefore describes the search's current empirical
knowledge, but it does not by itself determine or enumerate the full
hypothesis space.

Using this surrogate-relative notion of knowledge, we partition $\X$ into four
epistemic regimes, as illustrated in
Figure~\ref{fig:knowledge-regimes}.

\begin{itemize}[leftmargin=1.6em,itemsep=2pt,topsep=2pt]
    \item \emph{Known knowns} are evaluated designs whose rewards are
    sufficiently well understood. The surrogate predicts their values with
    low posterior uncertainty $\sigma_t(x)$.

    \item \emph{Known unknowns} are designs that the search can formulate and
    the surrogate can model, but whose rewards remain unresolved. They are
    represented by high posterior uncertainty $\sigma_t(x)$ and can be
    resolved through further evaluation.

\item \emph{Unknown knowns} are designs or evaluation results that exist in
    an external source, such as a database, a previous experiment, or a
    parallel discovery process, but are absent from the current search
    context. They can become available only through an explicit retrieval
    operation.
    
    \item \emph{Unknown unknowns} are designs that lie beyond the search's
    current effective reach. This can occur for two reasons:
    (i)~$\X$ is too large, combinatorial, or effectively unbounded for the current
    search procedure to reach the design; or (ii)~the design lies outside the
    surrogate's modelling support, so that $(\mu_t(x),\sigma_t(x))$ does not
    constitute a meaningful calibrated prediction.
\end{itemize}

These regimes distinguish uncertainty, reachability, modelling, and
knowledge. A design does not become known merely because it belongs to a
mathematically defined space; a reachable design is not necessarily supported
by the surrogate; and information stored externally is not available unless
it is retrieved into the current context. All four regimes may therefore
persist throughout the search.

The \emph{search frontier} is the conceptual boundary of what the current
procedure can effectively formulate, reach, and model. Designs inside this
frontier are accessible to the search, although their rewards may remain
uncertain. Designs outside it have not yet entered the modelled search
process. This distinction gives rise to three operations closely related to \cite{wang2008infinitely}:

\begin{itemize}[leftmargin=1.6em,itemsep=2pt,topsep=2pt]
    \item \emph{Exploitation} acts on known knowns by selecting designs for
    which the surrogate predicts both high reward $\mu_t(x)$ and low
    uncertainty $\sigma_t(x)$.

    \item \emph{Exploration} resolves known unknowns by evaluating reachable
    but uncertain designs. The resulting observations reduce posterior
    uncertainty and may turn known unknowns into known knowns.

    \item \emph{Discovery} expands the search frontier by formulating
    hypotheses that were previously outside the search's effective reach.
    This brings unknown unknowns into the modelled hypothesis space, where they
    become known unknowns that can subsequently be evaluated.
\end{itemize}

Discovery is therefore not simply the evaluation of an untested design. An
untested design that already lies within the surrogate's modelling support is
a known unknown and belongs to exploration. Discovery instead changes what
the search can express, reach, or meaningfully model. Experimental evaluation
then determines the value of the newly surfaced hypothesis.

The three operations correspond to three transitions in
Figure~\ref{fig:knowledge-regimes}: discovery moves hypotheses from unknown
unknowns to known unknowns, exploration moves them from known unknowns to
known knowns, and exploitation operates within the known-known regime. The
unknown-known regime requires separate memory read and write operations. We
leave memory retrieval and federated discovery to future work and focus here
on discovery, exploration, and exploitation within a single sequential search
process.

This formulation is conceptual and does not prescribe a particular discovery
algorithm. The next section introduces the Large Discovery Model, which
instantiates these operations by combining a generative foundation model, the
probabilistic surrogate introduced above, and an acquisition-guided
inference-time search policy.


